% ============================================================
% 跨设备脑电数据融合技术综述
% Cross-Device EEG Data Fusion: A Comprehensive Survey
% ============================================================
\documentclass[12pt, a4paper]{article}

% ---------- 中文排版 ----------
\usepackage[UTF8, fontset=fandol]{ctex}

% ---------- 页面布局 ----------
\usepackage[top=2.5cm, bottom=2.5cm, left=3cm, right=3cm]{geometry}

% ---------- 数学与符号 ----------
\usepackage{amsmath, amssymb, amsfonts}

% ---------- 图表 ----------
\usepackage{graphicx}
\usepackage{booktabs}
\usepackage{tabularx}
\usepackage{longtable}
\usepackage{array}
\usepackage{multirow}
\usepackage{float}

% ---------- 超链接与引用 ----------
\usepackage[colorlinks=true, linkcolor=blue, citecolor=blue, urlcolor=blue]{hyperref}
\usepackage{natbib}
\bibliographystyle{unsrtnat}

% ---------- 其他 ----------
\usepackage{xcolor}
\usepackage{enumitem}
\usepackage{setspace}
\onehalfspacing

% ---------- 标题信息 ----------
\title{%
  \Large\textbf{跨设备脑电数据融合技术综述}\\[0.5em]
  \large Cross-Device EEG Data Fusion: A Comprehensive Survey
}
\author{%
  综述报告 · 2026年4月
}
\date{}

% ============================================================
\begin{document}
\maketitle
\tableofcontents
\newpage

% ============================================================
\begin{abstract}
\noindent
脑电图（EEG）作为一种高时间分辨率、低成本的神经成像技术，在脑机接口、认知神经科学和临床神经病学中具有举足轻重的地位。然而，长期以来，EEG数据的采集呈现高度碎片化的格局：不同研究机构使用来自不同厂商的设备，采用迥异的参考方式、采样率、电极布局与文件格式，导致各数据集之间难以直接比较和融合，形成严重的"数据孤岛"效应。这一局面严重制约了大规模脑电基础模型的训练与泛化能力。

本综述系统梳理了跨设备EEG数据融合领域的最新研究进展，涵盖以下核心主题：（1）主流EEG设备的硬件异构性分析，量化设备效应对信号方差的影响（约占总方差9\%）；（2）湿电极与干电极跨设备信号质量的实证比较；（3）球面样条插值、重参考统一化、鲁棒归一化等预处理与空间对齐策略；（4）BIDS-EEG规范与HED事件标注体系的元数据标准化方案；（5）ComBat统计协调、欧几里得对齐、黎曼流形几何对齐与对抗域适应等特征级融合方法；（6）从BENDR到LaBraM、BIOT、CBraMod、REVE、EEGPT的EEG基础模型发展脉络及其多设备自适应机制；（7）80+个公开数据集的系统梳理，涵盖多中心数据集、同步多设备对照数据集、按范式分类的主要数据集（运动想象、SSVEP、情绪识别、临床睡眠）、推荐跨设备融合配对方案、基础模型预训练语料来源及元数据质量现状；（8）五大关键研究空白与可行的未来技术路线。

总体判断：原始信号级跨设备融合可行性中等偏低，特征级融合可行性较高，模型级迁移与域自适应可行性最高。随着灵活分词机制与四维位置编码的引入，EEG基础模型已初步具备设备无关化表征学习能力，但仍需解决设备效应显式建模、traveling-subject基准数据集构建与跨设备评估标准化等核心问题。
\end{abstract}

% ============================================================
\section{引言}
\label{sec:intro}

脑电图（Electroencephalography, EEG）自20世纪30年代问世以来，始终是神经科学研究的核心工具之一。相较于功能性磁共振成像（fMRI）和正电子断层扫描（PET），EEG具有毫秒级时间分辨率、设备成本相对低廉、受试者移动限制少等独特优势，被广泛应用于癫痫监测、睡眠分期、认知负荷评估、情绪识别以及脑机接口（Brain-Computer Interface, BCI）系统的实时控制中\citep{Lotte2018reviewBCI}。

进入深度学习时代，语言模型和视觉模型的成功在很大程度上依赖于"互联网规模"的同质数据集——所有文本具有统一的字符编码，所有图像具有兼容的像素格式。然而，EEG数据的生态系统却呈现出截然不同的面貌：全球数百个实验室使用从BioSemi ActiveTwo、Brain Products actiCHamp到Emotiv EPOC、OpenBCI Cyton、Muse等数十种设备，在各异的参考方案、硬件滤波配置、电极数量（从4个到256个不等）和采样率（从128\,Hz到16384\,Hz不等）下独立采集数据，并以相互不兼容的私有格式存储。这种高度碎片化的现状导致了严重的"数据孤岛"效应：单个研究的典型样本量仅有数十名受试者，在特定设备上训练的模型在面对不同设备时性能骤降\citep{EEGFMreview2025}。

这一局面的代价是双重的。从科学层面看，我们无法确定深度学习模型学到的究竟是真实的神经生理规律，还是设备特有的"硬件指纹"。从工程层面看，面向临床部署的BCI系统若仅能在特定设备上运行，其推广价值将大打折扣。正因如此，将来自不同设备的EEG数据融合为统一的、可互操作的大规模数据集，不仅是构建通用脑电基础模型的先决条件，也是推动EEG技术从实验室迈向实际应用的关键一步。

从技术进展来看，该领域近年来迎来了多个重要突破：

\begin{enumerate}[leftmargin=2em]
  \item \textbf{统计协调方法的引入}：2024年，Jaramillo-Jimenez等人首次将ComBat方法系统性地应用于多站点EEG数据的批次效应校正\citep{Jaramillo2024ComBat}，而此类方法在fMRI领域的应用比EEG早了约八年\citep{Fortin2017DTI}；
  \item \textbf{EEG基础模型的涌现}：2021年至2025年间，BENDR\citep{Kostas2021BENDR}、LaBraM\citep{Jiang2024LaBraM}、BIOT\citep{Yang2023BIOT}、CBraMod\citep{Wang2025CBraMod}、REVE\citep{ElOuahidi2025REVE}等一系列在多设备、多数据集上预训练的大规模模型相继出现，通过灵活的分词机制天然支持异构电极配置；
  \item \textbf{同步对照实验的实证支撑}：D'Amico与de Sa（2026）的同时采集研究表明，OpenBCI Cyton与Brain Products BrainAmp在同一头皮位置的记录可达97\%--100\%的信号相似度\citep{DAmico2026}。
\end{enumerate}

本综述的目标是为研究者提供该领域从硬件基础到算法前沿的全景视图，并系统梳理当前主要研究空白与未来技术路线。文章结构如下：第\ref{sec:hardware}章分析设备硬件异构性；第\ref{sec:quality}章综述跨设备信号质量比较研究；第\ref{sec:preprocessing}章介绍预处理与空间对齐策略；第\ref{sec:metadata}章讨论元数据标准化体系；第\ref{sec:feature}章综述特征级融合与域适应方法；第\ref{sec:foundation}章介绍EEG基础模型及其多设备自适应机制；第\ref{sec:datasets}章梳理开放数据资源与评估基准；第\ref{sec:gaps}章分析研究空白与未来方向；第\ref{sec:conclusion}章给出结论。

% ============================================================
\section{EEG系统硬件异构性分析}
\label{sec:hardware}

要科学评估跨设备数据融合的可行性，首先必须深入理解阻碍融合的底层物理障碍。市场上的EEG设备在电极材质、参考机制、放大器设计、数模转换精度、硬件滤波策略以及数据格式等多个维度上存在根本性的差异，这些差异共同构成了跨设备融合的技术鸿沟。

\subsection{主流设备硬件规格对比}

表\ref{tab:devices}给出了当前神经科学界与工业界最具代表性设备的硬件技术规格对比。

\begin{table}[H]
\centering
\caption{主流EEG设备硬件规格对比}
\label{tab:devices}
\small
\begin{tabularx}{\textwidth}{p{3cm}XXXX}
\toprule
\textbf{设备} & \textbf{电极类型} & \textbf{参考机制} & \textbf{采样率/ADC} & \textbf{滤波特性} \\
\midrule
BioSemi ActiveTwo & 有源湿电极 & CMS/DRL主动闭环，无传统REF/GND & 最高16\,kHz / 24位 & 纯DC耦合，五阶CIC数字低通 \\
\addlinespace
EGI Geodesic (HydroCel) & 盐水海绵湿电极 & 头顶Cz作物理参考 & 250--1000\,Hz / 光纤隔离 & ADAPT技术，高输入阻抗，4\,kV隔离 \\
\addlinespace
Brain Products actiCHamp Plus & 有源/无源可选 & 独立REF+GND接口，支持超扫描 & 最高100\,kHz / 24位独立ADC & 极宽带宽(DC至7500\,Hz) \\
\addlinespace
Emotiv EPOC X & 半干聚合物电极 & 内置类CMS/DRL降噪 & 128/256\,Hz / 14位 & 0.2--45\,Hz带宽，AC耦合 \\
\addlinespace
OpenBCI Cyton & 模块化干/湿可选 & 自定义单端或差分 & 200--250\,Hz / 24位 & ADS1299芯片，开源底层访问 \\
\addlinespace
Muse S/2 & 银墨水扁平干电极 & 内部硬件参考 & 256\,Hz / 12位 & 高度集成ASIC，仅4--6通道 \\
\bottomrule
\end{tabularx}
\end{table}

\subsection{参考逻辑体系的根本异构性}

EEG记录的本质是两点间的电位差，因此参考电极的选择从根本上决定了整个头皮电位图谱的形态。各主流设备在参考方案上的差异形成了一种"南辕北辙"的局面：

\textbf{BioSemi CMS/DRL系统}：BioSemi彻底废弃了传统的参考和接地电极概念，转而采用共模感知（Common Mode Sense, CMS）有源电极与驱动右腿（Driven Right Leg, DRL）无源电极构成的闭环反馈系统\citep{BioSemiSpecs}。该系统将受试者的平均共模电压持续钳制在ADC内部参考电压附近，赋予前端信号超过80\,dB的共模抑制比（CMRR）。用户必须在后期分析阶段手动指定数学参考（如全脑平均参考CAR）。

\textbf{EGI Geodesic系统}：EGI系统以头顶Cz作为默认硬件参考点\citep{BIDSspecEEG}，所有通道记录的均是相对于Cz的电位差。这与BioSemi的录制逻辑存在根本性差异，直接融合两者数据会引入严重的系统性偏差。

\textbf{Brain Products系统}：actiCHamp Plus为每个放大器模块提供独立的REF和GND接口，特别考虑了多被试超扫描场景下的电气隔离需求\citep{BrainProductsSpecs}。

\textbf{消费级设备}：Emotiv和Muse等消费级设备的参考方案高度集成于硬件内部，用户缺乏底层控制权，也缺乏精确的参考位置信息\citep{EmotivSpecs}。

上述差异要求跨设备数据融合管线必须实施统一的重参考操作（如平均参考或REST变换），否则融合后的特征空间将陷入逻辑混乱\citep{Chella2016reference}。

\subsection{放大器链路差异：滤波、阻抗与噪声本底}

放大器的硬件设计差异决定了信号的本底频谱特性：

\textbf{DC耦合vs.\ AC耦合}：BioSemi采用纯DC耦合设计，能无失真地记录包括慢波（$<$0.1\,Hz）在内的全频段信号\citep{BioSemiSpecs}；而Emotiv EPOC X采用AC耦合，其0.2\,Hz高通截止频率永久性地丢失了DC漂移信息，这使得两类设备在低频频段的信号根本上不可比。

\textbf{输入阻抗与抗噪能力}：BioSemi的前端输入阻抗在50\,Hz时高达300\,M$\Omega$，能最大程度地免疫工频干扰；消费级干电极设备的阻抗远低于此，导致更高的白噪声水平和更强的环境干扰敏感性\citep{Ratti2017systems}。

\textbf{ADC量化精度}：从Muse的12位（动态范围约72\,dB）到BioSemi/Brain Products的24位（动态范围约144\,dB），量化精度相差约4096倍，直接影响信号的数值分布范围与噪声底\citep{Sensors2024spectral}。

\subsection{设备效应的量化：约占总信号方差的9\%}

对跨设备EEG效应最严谨的量化研究来自Ratti等人（2017）\citep{Ratti2017systems}。该研究在六种标准实验范式下对EEG数据方差进行了系统分解：\textbf{被试间差异在总方差中占主导地位（约32\%），设备/系统差异约占9\%，会话间变异仅约1\%}。当样本量约为16名被试时，设备方差已与单个被试的个体方差量级相当——这意味着任何跨系统汇集数据的研究都不可忽视设备效应。

进一步的发现是：研究级系统之间具有统计上的可互换性（在各范式间无显著均值差异），而消费级设备会引入系统性偏差，在至少一半的测试范式中产生显著差异值。这一结论为跨设备融合策略的分级设计提供了重要依据。

% ============================================================
\section{跨设备信号质量比较研究}
\label{sec:quality}

认识到硬件规格的差异之后，一个自然的问题是：这些差异在实际测量信号中如何体现？本章综述跨设备信号质量的实证比较研究，覆盖湿电极与干电极的对比、ERP成分一致性、频谱特性差异，以及消费级设备在机器学习任务中的解码性能。

\subsection{湿电极vs.\ 干电极：信号质量实证}

传统观点认为，依赖导电膏的湿电极系统（如Brain Products、BioSemi）是信号保真度的唯一黄金标准。然而，随着材料科学与主动放大技术的发展，这一观点已受到系统性挑战。

\textbf{ERP成分的跨设备一致性}：在失配负波（Mismatch Negativity, MMN）实验中，干电极与湿电极记录的MMN峰值幅度在统计学上并无显著差异（$p=0.911, d=0.02$），充分证明了干电极在捕捉关键认知事件方面的有效性\citep{dryWetMMN2024}。在更早期的研究中，消费级设备同样能可靠地记录P100、N100、P200、N200和P300等晚期认知成分\citep{Sabio2024scoping}。

\textbf{D'Amico \& de Sa（2026）同步对照研究}：这项最新研究为跨设备比较提供了方法论上最为严格的"地基级证据"\citep{DAmico2026}。研究者将OpenBCI Cyton与Brain Products BrainAmp在同一批受试者、同一批任务、同一套头皮电极位置下进行\textbf{同时记录}，比较P3b、ERN、N400等ERP成分。结果显示两套系统的记录高度相似，信号相似度达97\%--100\%，MAE约为1\,$\mu$V量级。作者强调二者"高度相似但并非完全等价"，意味着量化精度、采样率和残余偏差的影响不可忽视。

\subsection{频谱特性差异}

频谱分析揭示了设备间更为微妙的差异。Belo等人（2024）对消费级与研究级设备的频谱特征进行了系统比较\citep{Sensors2024spectral}：

\begin{itemize}
  \item PSBD Headband Pro在alpha频段与Brain Products的相关性可达$r=0.9$，属于较好的消费级设备；
  \item Muse的信号质量最差，与研究级设备的频谱对齐度极低；
  \item 干电极设备在Delta和Theta等低频频段的功率谱密度异常升高，这与DC耦合缺失以及接触阻抗较高引起的低频噪声放大直接相关。
\end{itemize}

\subsection{消费级设备在ML任务中的解码性能}

消费级设备产生的数据并非"低质"到无法利用，但其解码策略需针对性调整。

\textbf{情绪识别任务}：在DREAMER数据集（14导联Emotiv设备采集）上，直接应用为高密度高信噪比数据设计的EEGNet时，F1分数仅为0.567。而采用统计、频域和连通性特征的领域特定特征工程，配合随机森林分类器，F1分数可达0.945，提升幅度高达67\%（$p<0.000001$, Cohen's $d=3.863$）\citep{EmotionCrossDataset2025}。这表明针对低信噪比数据需要专门的特征工程方案，而不能直接照搬高密度设备上的端到端深度学习管线。

\textbf{运动想象（MI）任务}：在受控环境下，湿电极系统在主体依赖的MI分类准确率上通常稳定在70\%至96\%之间，而干电极系统在约15\%至22\%的方差波动下，平均识别准确率依然能突破71\%的实用阈值\citep{dryEEGmotorimagery2025}。

\textbf{预处理影响}：Kessler等人（2025）的研究表明，不同预处理路径会显著改变解码性能，某些去伪迹步骤甚至会降低分类表现\citep{Kessler2025preprocessing}。这强调了"预处理一致性"必须作为统一数据集的明确协议加以规定。

% ============================================================
\section{信号预处理与空间对齐策略}
\label{sec:preprocessing}

跨设备数据融合的第一道工程关卡，是在原始信号维度实施高度标准化的对齐操作。本章系统介绍球面样条插值、重参考统一化、鲁棒归一化、重采样与抗混叠滤波，以及基于CNN的通道上采样方法。

\subsection{球面样条插值（Spherical Spline Interpolation, SSI）}

当两个数据集的通道数量和电极位置不同时，最直接的需求是将一种通道配置映射到另一种配置。球面样条插值（SSI）是目前解决这一问题的标准方法\citep{Perrin1989spherical}。

SSI的核心思想是将真实的头皮抽象为一个理想的三维单位球面。对于给定的源电极位置$\mathbf{r}_i$和目标虚拟电极位置$\mathbf{r}_j$，算法通过求解基于球面谐波的插值矩阵，在未布置真实电极的虚拟坐标点上，根据周围电极的电位梯度平滑地插值出合成信号：

\begin{equation}
V(\mathbf{r}_j) = \sum_{i=1}^{N} c_i \cdot g(|\mathbf{r}_i - \mathbf{r}_j|) + d
\end{equation}

其中$g(\cdot)$为球面样条基函数，$c_i$和$d$通过最小二乘法求解。相较于最近邻插值或平面样条插值，SSI充分考虑了颅骨几何形状与电场弥散特性，在地形图保真度与相位一致性上具有明显优势\citep{Shin2021CNNupsampling}。

2021年提出的参考电极标准化插值技术（RESIT）在SSI基础上进一步改进，误差降低了2.39\%至33.5\%，为高密度到低密度电极布局的转换提供了更精确的工具\citep{RESIT2021}。

\subsection{基于CNN的通道上采样}

Shin等人（2021）提出了一种基于卷积神经网络的通道上采样方法\citep{Shin2021CNNupsampling}，该方法在5144小时、1385名被试数据上训练后，可将4或14通道的稀疏EEG重建为21通道的密集EEG，经认证神经电生理学家判断，其质量与真实记录无法区分。这为低通道消费级设备数据的空间增强提供了一条数据驱动的路径。

\subsection{重参考统一化}

由于不同设备采用不同的参考方案，跨设备融合前必须将所有数据统一到同一参考体系下：

\begin{itemize}
  \item \textbf{公共平均参考（CAR）}：将所有通道的均值作为新参考，适用于覆盖范围较均匀的高密度电极布局；
  \item \textbf{标准化参考变换（REST）}：通过球面前向模型将测量数据重参考到无穷远处的参考点，理论上提供了最为客观的参考方案；
  \item \textbf{链接乳突参考}：将双侧乳突电极的均值作为参考，是临床应用中的传统选择。
\end{itemize}

Chella等人（2016）的研究指出，参考方式会显著影响头皮EEG连通性估计与图论指标\citep{Chella2016reference}，因此跨设备比较中参考方案的统一是必要条件而非可选项。

\subsection{鲁棒归一化}

不同设备放大器的模数转换增益差异会导致记录到的微伏级电压在绝对数值上存在极大的分布偏差。传统的Z-score标准化（基于均值和标准差）在跨设备融合中极其脆弱，因为EEG信号容易受到瞬态肌电爆发、电极脱落或眨眼引发的极端异常幅值污染。

鲁棒标准化技术\citep{EEGnormalization2025}使用数据集的\textbf{中位数}进行中心化，并利用\textbf{四分位距（IQR）}进行尺度缩放：

\begin{equation}
\tilde{x} = \frac{x - \mathrm{median}(x)}{\mathrm{IQR}(x)}
\end{equation}

实验表明，这种策略能有效压缩由异常毛刺带来的数值失真，显著降低跨记录和跨设备的信号变异性，使不同设备产生的波形特征在送入下游模型前收敛到相似的数值流形空间中。

\subsection{重采样与抗混叠滤波}

跨设备融合往往需要将所有数据重采样到统一的采样率（通常选择各数据集最低采样率，或约200--256\,Hz的标准值）。重采样前必须施加截止频率为目标采样率一半的低通抗混叠滤波器。应优先使用零相位FIR滤波器（如Hamming窗）而非IIR滤波器，以保证各通道间的相位一致性\citep{Bigdely2015PREP}。

\subsection{预处理管线的重要性与规范化}

Kessler等人（2025）的系统性研究强调，预处理选择对解码性能的影响不亚于模型架构的选择\citep{Kessler2025preprocessing}。MNE-BIDS-Pipeline等开源工具链的出现，为多设备、多中心数据的自动化、可复现预处理提供了工程支撑。

% ============================================================
\section{元数据标准化：BIDS与HED生态系统}
\label{sec:metadata}

仅在数值和空间层面实现对齐还不够——融合后的多源数据集在语义和组织结构上也必须具备互操作性。本章介绍脑成像数据结构（BIDS-EEG）和层级事件描述符（HED）两大标准化体系，以及MNE-BIDS-Pipeline自动化工具链。

\subsection{BIDS-EEG规范}

BIDS-EEG是对脑成像数据结构（Brain Imaging Data Structure）的EEG扩展\citep{Pernet2019BIDS}，制定了一套机器可读的目录层级和文件命名准则。其核心贡献在于：

\begin{enumerate}[leftmargin=2em]
  \item \textbf{统一文件格式}：来自不同厂商的私有格式（BioSemi的\texttt{.bdf}、Brain Products的\texttt{.vhdr/.eeg}、EGI的\texttt{.mff}）均需转换为统一的开放格式（如European Data Format \texttt{.edf}或BrainVision格式）。

  \item \textbf{强制性元数据字段}：BIDS规定了一系列对跨设备研究至关重要的字段，包括：
    \begin{itemize}
      \item \texttt{Manufacturer}（设备厂商）
      \item \texttt{ManufacturersModelName}（型号）
      \item \texttt{SamplingFrequency}（采样率）
      \item \texttt{EEGReference}（参考方式）
      \item \texttt{HardwareFilters}（硬件滤波参数）
    \end{itemize}

  \item \textbf{标准化通道与电极文件}：\texttt{\_channels.tsv}明确每个通道的类型（EEG/EOG/ECG/其他）、采样率及单位；\texttt{\_electrodes.tsv}存储每个电极的三维坐标（通过3D数字化仪或标准模板映射获得）。这种结构化存储使自动化脚本能准确提取几何和电物理属性进行动态校正。
\end{enumerate}

目前，OpenNeuro平台已托管超过60个遵循BIDS-EEG标准的数据集\citep{Markiewicz2021OpenNeuro}，为自动化跨设备数据融合提供了现实基础。

\subsection{HED层级事件描述符}

不同实验室对于刺激事件的记录习惯千差万别——同样的"左手运动想象"事件，在A实验室可能标记为\texttt{S 1}，在B实验室可能标记为\texttt{Marker\_Left\_Hand}。这种语义不一致性构成了自动化跨数据集训练的致命壁垒。

层级事件描述符（Hierarchical Event Descriptors, HED）\citep{Niso2022HED}提供了一套高度结构化的机器可读词汇库，能够精确描述实验中的事件类型和上下文。例如，左手运动想象可以用以下HED标注唯一地表达：

\begin{center}
\texttt{Action/Imagine/Upper-extremity/Hand, Direction/Left}
\end{center}

通过为BIDS数据集的\texttt{\_events.tsv}文件引入HED映射字典，来自不同实验室、不同设备的异构事件被归一化到同一概念空间，极大提升了数据集间交叉验证的潜力。

\subsection{MNE-BIDS-Pipeline自动化工具链}

MNE-BIDS-Pipeline将上述规范落地为可直接使用的自动化预处理管线。通过一个简单的文本配置文件，该工具链可对包含数百名被试、混杂多种设备格式的庞大BIDS数据集执行：滤波重采样→球面空间重映射→ICA伪迹剔除→特征矩阵提取的全流程并行化批处理。这标志着在数据工程层面，将不同设备采集的海量脑电信号整编为标准可用数据集已具备坚实的工具支撑。

% ============================================================
\section{特征级融合与域适应方法}
\label{sec:feature}

即便预处理阶段穷尽了插值、重构与标准化策略，不同设备残留的硬件频响差异、局部电极微偏以及不可观测的环境噪声仍会使高级特征产生不可忽视的分布偏移（Distribution Shift）。本章综述统计协调方法、欧几里得对齐、黎曼流形对齐、对抗域适应以及跨设备迁移学习四大类特征级融合方法。

\subsection{统计协调：ComBat及其变体}

\subsubsection{方法起源与EEG应用}

ComBat方法最初由Johnson等人（2007）在基因组学领域提出\citep{Johnson2007ComBat}，随后被Fortin等人引入神经影像领域用于协调多站点扩散张量成像和皮层厚度测量数据\citep{Fortin2017DTI, Fortin2018cortical}。然而，ComBat在EEG领域的首次系统应用直到2024年才由Jaramillo-Jimenez等人实现\citep{Jaramillo2024ComBat}——较fMRI领域滞后约八年。

ComBat的核心思想是通过加性和乘性混合效应模型校正批次效应：

\begin{equation}
Y_{ijg} = \alpha_g + X_{ij}\beta_g + \gamma_{ig} + \delta_{ig}\varepsilon_{ijg}
\end{equation}

其中$Y_{ijg}$为特征值，$\gamma_{ig}$和$\delta_{ig}$分别为加性和乘性批次效应，通过经验贝叶斯方法估计。

\subsubsection{比较研究结果}

Jaramillo-Jimenez等人（2024）系统比较了neuroCombat、neuroHarmonize、OPNested-GMM和HarmonizR四种变体在5个独立数据集（$n=374$）的汇集静息态EEG上的表现\citep{Jaramillo2024ComBat}。结果表明，\textbf{HarmonizR和OPNested-GMM ComBat在有效减少传感器空间频谱特征批次效应的同时，最好地保留了与年龄相关的生物学关联}。

\subsubsection{HarMNqEEG：黎曼几何协调}

具有里程碑意义的HarMNqEEG项目（Li、Valdes-Sosa等，2022）基于\textbf{12种不同EEG设备、9个国家、1564名被试}的数据，在交叉功率谱张量的黎曼几何框架下构建了全生命周期标准化定量EEG\citep{Li2022HarMNqEEG}。协调后的黎曼z分数提高了脑发育障碍的诊断准确率，并建立了迄今最全面的多设备EEG常模数据库。Engemann等人（2025）进一步证明，基于协方差技术的M/EEG数据协调能显著增强预测回归建模的能力\citep{HarmonizeCovariance2025}。

\subsection{欧几里得对齐（Euclidean Alignment, EA）}

欧几里得对齐（EA）由华中科技大学吴冬蕊团队（2020）提出\citep{He2020EA}，至今仍是最广泛采用的实用域适应方法之一。EA的核心操作是将不同被试或设备的EEG协方差矩阵$\bar{R}$对齐到公共参考空间：

\begin{equation}
\tilde{X}_i = \bar{R}_i^{-1/2} X_i
\end{equation}

其中$\bar{R}_i$是第$i$个域（被试或设备）的平均协方差矩阵，$\tilde{X}_i$是对齐后的信号。这一简洁的操作显著减少了跨域协方差分布的差异，成为NeurIPS 2021 BEETL（脑电迁移学习基准）竞赛获胜方案的核心组件\citep{RevisitEA2025}。

\subsection{黎曼流形几何对齐}

EEG信号的空间协方差矩阵驻留在对称正定（SPD）矩阵构成的黎曼流形上，而非平坦的欧几里得空间\citep{Barachant2013Riemannian}。传统欧氏距离算法在处理这类数据时会产生几何畸变。前沿的黎曼对齐策略\citep{RiemannianReview2024}设计了三个递进变换：

\begin{enumerate}[leftmargin=2em]
  \item \textbf{重新居中（Re-centering）}：在黎曼流形上计算各数据集的Fréchet均值$G_i$，将所有源数据集平移至流形上的统一参考点；
  \item \textbf{离散度均衡（Dispersion Equalization）}：缩放各数据集在质心周围的散布半径，使不同放大器噪声本底引起的分布差异得到补偿；
  \item \textbf{旋转校正（Rotation Correction）}：通过正切空间投影，应用相关对齐（CORAL）或普氏分析（Procrustes）修正源域与目标域特征轴间的旋转错位。
\end{enumerate}

基于此框架构建的BARN-DA模型通过最小化黎曼最大均值差异（R-MMD）损失，在跨会话与跨被试分类精度上取得了突破性进展\citep{BARNDA2025}。

\subsection{对抗域适应（DANN及衍生方法）}

领域对抗神经网络（DANN）\citep{Ganin2016DANN}通过引入梯度反转层（Gradient Reversal Layer），在优化任务分类精度的同时最大化域分类误差，迫使特征提取网络学习域不变表征。这一思想已被广泛应用于跨被试和跨设备的情绪识别、认知负荷监测等场景\citep{crossSubjectMI2025}。

\subsection{专门跨设备迁移学习方法}

表\ref{tab:transfer}总结了专门针对跨设备场景设计的迁移学习方法。

\begin{table}[H]
\centering
\caption{跨设备EEG迁移学习主要方法对比}
\label{tab:transfer}
\small
\begin{tabularx}{\textwidth}{lXXX}
\toprule
\textbf{方法} & \textbf{核心策略} & \textbf{应用场景} & \textbf{主要结论} \\
\midrule
AwAR \citep{Wu2017AwAR} & 主动加权自适应正则化 & 跨头戴设备ERP & 60--65个新设备样本即可恢复基线性能 \\
\addlinespace
ALPHA \citep{Liu2022ALPHA} & 子空间对齐+统计量池化 & 湿/干电极SSVEP-BCI & 部分设置下可接近免重校准 \\
\addlinespace
SDDA \citep{Liu2025SDDA} & 空间蒸馏+分布对齐 & 异构通道布局 & 跨异构设备仍可获得稳定收益 \\
\addlinespace
CDRC \citep{Zhang2025CDRC} & 对比预训练跨设备表征一致性 & 跨设备EEG识别 & 自监督学习降低标注依赖 \\
\bottomrule
\end{tabularx}
\end{table}

\textbf{AwAR}（Wu等，2017）\citep{Wu2017AwAR}直接面向"更换EEG头戴设备"后的迁移问题，证明旧设备上的有标注数据可帮助新设备快速校准，显著减少新设备标注样本需求。

\textbf{SDDA}（Liu等，2025）\citep{Liu2025SDDA}针对通道数量和布局不一致的异构设备，通过空间蒸馏将高维设备的额外通道信息迁移至低维或不同布局设备，结合输入层、特征层和输出层的分布对齐获得稳定收益。

\textbf{CDRC}（Zhang等，2025）\citep{Zhang2025CDRC}采用双分支Transformer嵌入的自监督对比预训练，直接以跨设备表征一致性为学习目标，在无需大量标注数据的情况下实现跨设备泛化。

% ============================================================
\section{面向设备无关处理的EEG基础模型}
\label{sec:foundation}

如果说域适应技术是在算法后端"被动修补"设备差异，那么近年涌现的大规模EEG基础模型则试图从架构设计的初始阶段就从根本上解决跨设备融合的障碍。本章梳理从BENDR到REVE的基础模型发展脉络，重点分析各模型的多设备自适应机制，并讨论缩放定律偏离与评估基准问题。

\subsection{发展脉络}

表\ref{tab:foundation}给出了主要EEG基础模型的对比。

\begin{table}[H]
\centering
\caption{主要EEG基础模型对比}
\label{tab:foundation}
\small
\begin{tabularx}{\textwidth}{lllXX}
\toprule
\textbf{模型} & \textbf{年份/会议} & \textbf{预训练规模} & \textbf{跨设备策略} & \textbf{机构} \\
\midrule
BENDR \citep{Kostas2021BENDR} & 2021, Frontiers & $\sim$1500小时 & 固定19通道输入 & 多伦多/Vector \\
\addlinespace
BIOT \citep{Yang2023BIOT} & 2023, NeurIPS & 多数据集 & 通道级分词，200\,Hz重采样 & UIUC/哈佛 \\
\addlinespace
LaBraM \citep{Jiang2024LaBraM} & 2024, ICLR Spotlight & $\sim$2500小时，$\sim$20数据集 & 通道-分片分词 & 上海交通大学 \\
\addlinespace
EEGPT \citep{EEGPT2024} & 2024, NeurIPS & 多数据集 & 掩码双重自监督 & 多机构 \\
\addlinespace
CBraMod \citep{Wang2025CBraMod} & 2025, ICLR & $\sim$9000小时 & 条件位置编码 & 浙江大学 \\
\addlinespace
REVE \citep{ElOuahidi2025REVE} & 2025, NeurIPS & $\sim$60000小时，25000名被试 & 四维傅里叶位置编码 & IMT Atlantique/蒙特利尔 \\
\bottomrule
\end{tabularx}
\end{table}

\subsection{通道-分片分词与位置编码机制}

传统卷积网络（如EEGNet\citep{Lawhern2018EEGNet}、ShallowConvNet\citep{Schirrmeister2017EEGNet}）要求固定的输入张量维度（通道数$\times$时间步数），面对异构电极配置时无能为力。

\textbf{LaBraM}（Jiang等，ICLR 2024 Spotlight）\citep{Jiang2024LaBraM}引入\textbf{通道-分片分词（channel-patch tokenization）}——将每个EEG通道独立分割为时间片段作为独立token，天然适应不同电极数量。该模型在BioSemi、BrainAmp、BCI2000、g.tec等约20个数据集上预训练约2500小时，参数规模从580万到3.69亿，在异常检测、事件分类、情绪识别和步态预测等任务上取得了领先结果。

\textbf{BIOT}（Yang等，NeurIPS 2023）\citep{Yang2023BIOT}采用通道级分词策略：将所有信号重采样至200\,Hz，每个通道编码为独立token并附加可学习的通道嵌入，天然支持不匹配的通道数、可变长度和缺失值。这使BIOT能够在4通道的Muse数据和128通道的EGI数据上无缝运行，而无需强行截断或补零。

\subsection{跨尺度时空分词：CSBrain}

CSBrain\citep{CSBrain2025}采用\textbf{跨尺度时空分词（Cross-Scale Spatiotemporal Tokenization, CST）}与结构化稀疏注意力（SSA）架构，依据大脑解剖学分区（额叶、顶叶、枕叶等）对通道信号进行动态聚合，使模型能够灵活应对不断变化的电极组合和空间分辨率。

\subsection{REVE：四维位置编码——最先进的跨设备方案}

REVE（El Ouahidi等，NeurIPS 2025）\citep{ElOuahidi2025REVE}代表了\textbf{迄今最先进的跨设备EEG处理方案}。其核心创新是\textbf{四维位置编码}——将每个电极映射到其三维空间坐标加时间位置，使用傅里叶嵌入编码：

\begin{equation}
\mathbf{e}(\mathbf{r}, t) = \left[\cos(2\pi \mathbf{W} [\mathbf{r}; t]), \sin(2\pi \mathbf{W} [\mathbf{r}; t])\right]
\end{equation}

其中$\mathbf{r} \in \mathbb{R}^3$为电极三维坐标，$t$为时间位置，$\mathbf{W}$为学习得到的频率矩阵。这种设计使得系统通过查询标准电极坐标数据库即可处理任意电极排列，实现在不重新训练的情况下的真正跨配置迁移。

REVE在约25000名被试、60000小时数据上预训练，平均性能超越先前模型约2.5\%，线性探测场景下提升可达17\%。

\subsection{其他近期模型}

\textbf{EEGPT}（NeurIPS 2024）\citep{EEGPT2024}采用基于掩码的双重自监督学习，同时对EEG信号的幅值和频率成分进行重建。\textbf{大认知模型（LCM）}\citep{LCM2025}引入可学习的通道映射矩阵，在预训练阶段显式学习不同电极配置间的映射关系。\textbf{BrainOmni}\citep{BrainOmni2025}作为首个EEG/MEG统一基础模型，通过模态嵌入实现跨模态预训练。\textbf{EEG-X}\citep{EEGX2025}被描述为"设备无关且噪声鲁棒"，通过基于位置的通道嵌入机制处理任意电极配置。

\subsection{缩放定律偏离与评估基准}

随着大规模融合数据集预训练的尝试，研究界发现了EEG领域异于文本和视觉任务的独特规律。EEG-FM-Bench\citep{EEGFMBench2025}在覆盖认知负荷、ERP、情绪检测、癫痫筛查等多种范式的综合测评中揭示了一个颇具警示意义的结论：

\textbf{在当前EEG数据规模和预训练目标下，简单增加模型参数量和数据量并不必然带来下游泛化性能的线性飞跃}——即"缩放定律"在极低信噪比的EEG任务中发生了偏离\citep{EEGFMworth2024}。原因在于：

\begin{itemize}
  \item EEG信号的高度非平稳性与严重环境噪声导致预训练重构目标的梯度与下游判别任务所需梯度产生冲突；
  \item 通过冻结预训练主干仅训练最后一层分类器的线性探测方法在跨设备验证中表现较弱；
  \item 完整的全参数微调对于发挥预训练泛化表征的迁移优势仍是必要的。
\end{itemize}

这些发现揭示，跨设备数据的物理融合与架构适应只是迈出了第一步，还需要在神经生理学底层规律和更贴合EEG特性的自监督目标上寻求突破。

值得注意的是，不同基础模型之间预训练数据的组成差异极大（详见第\ref{sec:pretrain_corpora}节）：BENDR\citep{Kostas2021BENDR}、CBraMod\citep{Wang2025CBraMod}和REVE\citep{ElOuahidi2025REVE}主要依赖TUH EEG语料库的多代Natus/Nicolet硬件数据；而LaBraM\citep{Jiang2024LaBraM}组合了约20个跨厂商数据集（BrainAmp、NeuroScan、BioSemi、BCI2000等）。这种预训练语料构成的差异直接影响下游跨设备泛化能力，但目前尚无研究系统性地控制预训练数据组成来量化其对设备泛化的独立贡献——这是基础模型评估体系中亟待填补的重要空白。

% ============================================================
\section{开放数据资源与评估基准}
\label{sec:datasets}

推动跨设备EEG融合研究的基础是高质量的开放数据资源和标准化评估基准。目前公开可用的EEG数据集已超过80个，涵盖10余种范式类型，使用了20余种不同的采集设备；仅MOABB框架就索引了来自3,421名被试的148个标准化BCI数据集\citep{Chevallier2024MOABB}。本章系统梳理多中心、多设备及范式特定数据集，并分析其跨设备融合潜力。

\subsection{多中心与多设备静息态数据集}

\textbf{HarMNqEEG}：Harmonized-Multinational qEEG Norms项目\citep{Li2022HarMNqEEG}是迄今唯一专为多设备协调而构建的数据集，汇集9个国家15项研究的1,564名健康受试者静息态数据，采用\textbf{12种不同的EEG放大器}。该数据集仅共享交叉功率谱张量而非原始EEG，限制了时域融合应用，但已验证其黎曼几何协调框架可提升脑发育障碍的诊断准确率。

\textbf{I-CARE}：国际心脏骤停研究联盟数据库\citep{ICARE2023}来自美国和欧洲7家不同医院，包含1,020名心脏骤停后昏迷患者超过5万小时的连续EEG数据，是多中心临床跨站点融合研究的代表性资源。

\textbf{HBN-EEG}：Child Mind Institute健康大脑网络数据集\citep{Alexander2017HBN}，包含超过3,000名儿童到年轻成人（5--21岁）的128通道EGI HydroCel EEG数据，被广泛用于EEG基础模型的训练与评估。

\textbf{LEMON}：莱比锡MPI大脑网络数据集\citep{Babayan2019LEMON}，使用Brain Products BrainAmp和actiCHamp采集228名健康成人（20--77岁）的62通道、2500\,Hz静息态EEG，以BIDS格式发布于OpenNeuro（ds000221），元数据完整性高，适合作为跨设备研究的参考域。

\textbf{TD-BRAIN}：Compumedics Neuroscan录制的1,274名精神科患者26通道静息态EEG\citep{vanDijk2022TDBRAIN}，发布于Synapse平台，提供了大规模临床群体的跨设备研究基础。

\subsection{同步多设备对照数据集}

同时或顺序使用多种设备采集同一被试的数据集，是验证跨设备对齐算法最直接有效的资源，但也是当前最为稀缺的数据类型。

\textbf{D'Amico \& de Sa（2026）}：目前方法论最严格的跨设备同步对照研究\citep{DAmico2026}，将OpenBCI Cyton与Brain Products BrainAmp在同一批受试者的相同头皮位置\textbf{同时}录制，比较P3b、ERN、N400等ERP成分，信号相似度达97\%--100\%。

\textbf{Lee等（2026）消费级-研究级对比数据集}\citep{Lee2026dataset}：30名被试\textbf{顺序佩戴} BrainLink Pro（1通道）、NeuroNicle FX2（2通道）、MindWave Mobile 2（1通道）、Muse 2（4通道）和DSI-24（21通道）5种设备，执行眨眼、咬牙、头动与静息态等任务，EDF+MAT双格式发布。

\textbf{Williams等（2020）OSF同步双设备数据集}\citep{Williams2020PeerJ}：在20名成人上\textbf{同步}录制Emotiv EPOC Flex（16通道，129\,Hz）与Neuroscan SynAmps2（16通道，1000\,Hz），涵盖MMN、P300、N170、SSVEP和静息态5种范式，原始数据发布于osf.io/zj3f5/，是目前对开发跨设备对齐算法最直接可用的多范式同步资源。

\textbf{Mobile BCI scalp+ear数据集}\citep{MobileBCI2021}：24名被试在步行和跑步时\textbf{同时}使用Brain Products BrainAmp（头皮，32通道，500\,Hz）和SMARTING cEEGrid（耳周，14通道，500\,Hz）采集ERP与SSVEP任务，为跨设备与跨记录位置的双重域偏移研究提供了独特资源。

\textbf{MoreGrasp湿/水基/干电极对比}\citep{Schwarz2020MoreGrasp}：BNCI Horizon 2020数据集，45名参与者（每类电极系统各15人）使用凝胶电极、水基电极和干电极系统执行相同的抓取运动想象任务，为电极类型对MI解码性能的系统性影响提供了实证数据。

\textbf{MAMEM1/MAMEM3配对}：MAMEM SSVEP数据库\citep{MAMEM2016}包含两个使用相同刺激协议（5个SSVEP频率，6.66--12\,Hz）但不同设备录制的配对子集：MAMEM1使用EGI GES 300（256通道，250\,Hz）；MAMEM3使用Emotiv EPOC（14通道，128\,Hz）。同一批11名被试在两种设备上均参与实验，提供了研究级vs.消费级SSVEP对比的纯粹基准。

\textbf{Badcock等（2015）儿童双设备验证}\citep{Badcock2015}：在19名儿童上同步录制Emotiv EPOC与Neuroscan（听觉oddball范式），为消费级设备在儿童群体中的ERP有效性提供了早期实证。

\subsection{按范式分类的主要数据集}

\subsubsection{运动想象数据集}

运动想象是跨设备融合研究最丰富的范式，拥有25+个开源数据集，横跨至少9种不同EEG系统，Brain Products BrainAmp系列出现在约9个数据集中，g.tec系统约7个，Neuroscan SynAmps2约5个。表\ref{tab:mi_datasets}列出最具代表性的MI数据集。

\begin{table}[H]
\centering
\caption{主要运动想象（MI）EEG数据集对比}
\label{tab:mi_datasets}
\small
\begin{tabularx}{\textwidth}{p{3.0cm}p{2.8cm}ccp{0.6cm}XX}
\toprule
\textbf{数据集} & \textbf{设备} & \textbf{通道} & \textbf{采样率} & \textbf{被试} & \textbf{类别} & \textbf{格式} \\
\midrule
BCI-IV-2a & BrainAmp MR plus & 22 & 250\,Hz & 9 & 4类 & GDF \\
\addlinespace
PhysioNet EEGMMIDB \citep{PhysioNetMotor} & BCI2000系统 & 64 & 160\,Hz & 109 & 4类(MI+ME) & EDF+ \\
\addlinespace
Lee2019\_MI \citep{Lee2019OpenBMI} & Brain Products BrainAmp & 62 & 1000\,Hz & 54 & 2类 & MAT \\
\addlinespace
Cho2017 \citep{Cho2017MI} & BioSemi ActiveTwo & 64 & 512\,Hz & 52 & 2类 & MAT \\
\addlinespace
Dreyer2023 \citep{Dreyer2023MI} & g.tec g.USBamp & 27 & 512\,Hz & 87 & 2类 & GDF \\
\addlinespace
Schirrmeister2017 \citep{Schirrmeister2017EEGNet} & NeurOne (Mega Electronics) & 128 & 500\,Hz & 14 & 4类 & EDF \\
\addlinespace
MoreGrasp \citep{Schwarz2020MoreGrasp} & 凝胶/水基/干电极 & 32 & 512\,Hz & 45 & 抓取MI & GDF \\
\bottomrule
\end{tabularx}
\end{table}

\subsubsection{SSVEP数据集}

SSVEP数据集横跨研究级到消费级设备，涵盖5家不同设备厂商。表\ref{tab:ssvep_datasets}列出主要数据集，MAMEM1/MAMEM3配对（EGI 256通道 vs. Emotiv EPOC 14通道）尤其适合研究设备质量对频率锁定信号的影响。

\begin{table}[H]
\centering
\caption{主要SSVEP EEG数据集对比}
\label{tab:ssvep_datasets}
\small
\begin{tabularx}{\textwidth}{p{3.0cm}p{2.8cm}ccp{0.6cm}X}
\toprule
\textbf{数据集} & \textbf{设备} & \textbf{通道} & \textbf{采样率} & \textbf{被试} & \textbf{刺激频率} \\
\midrule
Benchmark SSVEP \citep{WangBenchmarkSSVEP2016} & Neuroscan SynAmps2 & 64 & 250\,Hz & 35 & 40个(8--15.8\,Hz) \\
\addlinespace
BETA \citep{BETA2020} & Neuroscan SynAmps2 & 64 & 1000\,Hz & 70 & 40个(8--15.8\,Hz) \\
\addlinespace
Lee2019\_SSVEP \citep{Lee2019OpenBMI} & BrainAmp & 62 & 1000\,Hz & 54 & 4个(5.45--12\,Hz) \\
\addlinespace
MAMEM1 \citep{MAMEM2016} & EGI GES 300（研究级） & 256 & 250\,Hz & 11 & 5个(6.66--12\,Hz) \\
\addlinespace
MAMEM3 \citep{MAMEM2016} & Emotiv EPOC（消费级） & 14 & 128\,Hz & 11 & 5个(6.66--12\,Hz) \\
\addlinespace
Nakanishi2015 & BioSemi ActiveTwo & 8 & 256\,Hz & 10 & 12个 \\
\bottomrule
\end{tabularx}
\end{table}

\subsubsection{情绪识别数据集——三层设备梯度}

情绪识别数据集清晰地分为三个设备层级：研究级62通道（SEED系列，NeuroScan）、研究级32通道（DEAP/MAHNOB-HCI，BioSemi）和消费级14通道（DREAMER/AMIGOS，Emotiv EPOC）。这种三层梯度天然适合研究信号质量如何随设备类别衰减，同时保持范式一致性（表\ref{tab:emotion_datasets}）。

\begin{table}[H]
\centering
\caption{情绪识别EEG数据集设备层级对比}
\label{tab:emotion_datasets}
\small
\begin{tabularx}{\textwidth}{p{2.5cm}p{3.2cm}ccp{0.6cm}XX}
\toprule
\textbf{数据集} & \textbf{设备（层级）} & \textbf{通道} & \textbf{采样率} & \textbf{被试} & \textbf{情绪类别} \\
\midrule
SEED \citep{Zheng2015SEED} & ESI NeuroScan（研究级Ⅰ） & 62 & 200\,Hz & 15 & 3类（正/负/中） \\
\addlinespace
SEED-IV \citep{SEEDIV2018} & ESI NeuroScan（研究级Ⅰ） & 62 & 200\,Hz & 15 & 4类 \\
\addlinespace
DEAP \citep{Koelstra2012DEAP} & BioSemi ActiveTwo（研究级Ⅱ） & 32 & 128\,Hz & 32 & 效价/唤醒 \\
\addlinespace
MAHNOB-HCI & BioSemi ActiveTwo（研究级Ⅱ） & 32 & 256\,Hz & 27 & 效价/唤醒 \\
\addlinespace
DREAMER \citep{DREAMER2018} & Emotiv EPOC（消费级） & 14 & 128\,Hz & 23 & 效价/唤醒/支配度 \\
\addlinespace
AMIGOS & Emotiv EPOC（消费级） & 14 & 128\,Hz & 40 & 效价/唤醒 \\
\bottomrule
\end{tabularx}
\end{table}

\subsubsection{临床与睡眠数据集}

\textbf{TUH EEG语料库}（TUEG）\citep{Obeid2016TUH}是世界最大的开源临床EEG语料库，包含14,987名被试的69,652条记录、总计27,062小时，使用从2002年至今的多代Natus/Nicolet硬件录制，天然形成单一数据库内的多设备语料。其子集TUAB（异常检测）和TUSZ（癫痫事件分类）是EEG基础模型预训练的核心数据源。

睡眠数据集：Sleep-EDF Expanded（197条PSG记录，PhysioNet）、SHHS（约5,804条记录，Compumedics P-Series）\citep{SHHS2000}、MASS（200名被试，5家医院不同临床PSG系统）和ISRUC-Sleep（118名被试，3个亚组）\citep{ISRUC2016}共同构成睡眠跨设备研究的主要资源库。MASS的5个子集（SS1--SS5）在不同医院使用不同临床PSG系统录制，特别适合跨设备睡眠分期研究。

\subsubsection{静息态数据集}

LEMON\citep{Babayan2019LEMON}（Brain Products，228人）、TD-BRAIN\citep{vanDijk2022TDBRAIN}（Neuroscan，1,274人）和HBN-EEG\citep{Alexander2017HBN}（EGI，3,000+人）覆盖了三个主流研究级设备厂商，联合使用可构成多中心跨设备静息态研究矩阵。HarMNqEEG\citep{Li2022HarMNqEEG}（12种设备，9国，1,564人）是唯一专门设计的多设备协调数据集，提供了理论上最纯粹的设备效应量化基准。

\subsection{推荐跨设备融合数据集配对方案}
\label{sec:dataset_pairing}

本节基于"相同范式+不同设备"原则，梳理当前文献中最具可操作性的跨设备融合数据集配对方案，可作为第\ref{sec:gaps}章研究空白中标准化评估基准建设的可行起点。

\textbf{运动想象（四厂商，302人）}：PhysioNet EEGMMIDB\citep{PhysioNetMotor}（BCI2000，64通道，160\,Hz，109人）+ Lee2019\_MI\citep{Lee2019OpenBMI}（BrainAmp，62通道，1000\,Hz，54人）+ Cho2017\citep{Cho2017MI}（BioSemi，64通道，512\,Hz，52人）+ Dreyer2023\citep{Dreyer2023MI}（g.tec，27通道，512\,Hz，87人）。四个数据集共302名被试，横跨4家厂商，均采用左/右手运动想象范式，通道对齐可使用10-20公共子集（至少C3、C4、Cz），重采样至160\,Hz公共频率。

\textbf{SSVEP（研究级vs.消费级）}：MAMEM1（EGI 256通道）vs. MAMEM3（Emotiv 14通道）\citep{MAMEM2016}——相同11名被试、相同5个刺激频率，目前唯一已知的同被试研究级/消费级SSVEP配对。补充：Benchmark SSVEP + Lee2019\_SSVEP（同为研究级，不同厂商）。

\textbf{P300（三厂商）}：Lee2019\_ERP\citep{Lee2019OpenBMI}（BrainAmp，62通道，54人）+ EPFL P300（BioSemi，32通道，9人）+ BI2014a（g.tec干电极，16通道，71人）。

\textbf{情绪识别（三层设备梯度）}：SEED\citep{Zheng2015SEED}（NeuroScan，62通道）+ DEAP\citep{Koelstra2012DEAP}（BioSemi，32通道）+ DREAMER\citep{DREAMER2018}（Emotiv，14通道）——相同效价/唤醒标注框架，三级设备质量递进，是消费级-研究级桥梁研究的天然试验田。

\textbf{湿电极vs.干电极}：BI2012/BI2013a（g.tec湿电极，16通道）vs. BI2014a（g.tec干电极，16通道）——相同放大器，不同传感器类型。MoreGrasp\citep{Schwarz2020MoreGrasp}进一步扩展为凝胶/水基/干电极三向对比。

\textbf{头皮vs.耳周}：Mobile BCI数据集\citep{MobileBCI2021}（BrainAmp头皮32通道 + SMARTING cEEGrid耳周14通道，同步采集）是目前跨记录位置融合研究的唯一已知公开资源。

\subsection{基础模型预训练语料来源}
\label{sec:pretrain_corpora}

EEG基础模型的大规模预训练已验证了多数据集聚合的可行性。理解各模型的预训练数据组成，揭示了哪些数据集组合已被实际验证为可互操作的跨设备语料（表\ref{tab:pretrain}）。

\begin{table}[H]
\centering
\caption{主要EEG基础模型预训练数据组成}
\label{tab:pretrain}
\small
\begin{tabularx}{\textwidth}{p{1.8cm}p{0.8cm}Xp{1.8cm}X}
\toprule
\textbf{模型} & \textbf{年份} & \textbf{主要预训练数据源} & \textbf{规模} & \textbf{跨设备覆盖} \\
\midrule
BENDR \citep{Kostas2021BENDR} & 2021 & TUH EEG Corpus（完整） & $\sim$1,500小时 & Natus/Nicolet多代设备 \\
\addlinespace
BIOT \citep{Yang2023BIOT} & 2023 & SHHS + PREST + CHB-MIT + TUAB + TUEV & $\sim$1,000万样本 & Compumedics + Natus + 临床系统 \\
\addlinespace
LaBraM \citep{Jiang2024LaBraM} & 2024 & $\sim$20个多样数据集（含TUAB、SEED、DEAP、PhysioNet MMI、Sleep-EDF等） & $\sim$2,500小时 & BrainAmp + NeuroScan + BioSemi + BCI2000 \\
\addlinespace
CBraMod \citep{Wang2025CBraMod} & 2025 & TUEG（过滤后） & $\sim$9,000小时 & Natus/Nicolet多代 \\
\addlinespace
REVE \citep{ElOuahidi2025REVE} & 2025 & TUEG + OpenNeuro + HBN & $\sim$60,000小时，25,000人 & Natus + EGI + Brain Products \\
\addlinespace
EEGPT \citep{EEGPT2024} & 2024 & 混合多任务数据集 & $\sim$246小时 & 多厂商 \\
\bottomrule
\end{tabularx}
\end{table}

LaBraM的预训练方案\citep{Jiang2024LaBraM}对跨设备融合最具参考价值：它成功组合了约20个数据集，包括来自BrainAmp、NeuroScan、BioSemi和BCI2000等截然不同设备的数据，证明了\textbf{在适当通道对齐和归一化下，来自不同设备的EEG数据可以被有意义地聚合用于预训练}。各模型预训练数据的进一步差异分析见第\ref{sec:foundation}章；第\ref{sec:gaps}章将讨论这一多样性组合与当前基础模型未显式建模设备效应之间的深层矛盾。

\subsection{数据平台与基准评测工具}

\textbf{MOABB}（Mother of All BCI Benchmarks）\citep{Chevallier2024MOABB}：目前\textbf{最全面的EEG-BCI可复现性基准平台}，索引了来自3,421名被试的\textbf{148个标准化数据集}，支持Python编程接口自动下载和预处理，涵盖MI、P300、SSVEP、c-VEP和静息态范式。尽管MOABB涵盖了使用不同硬件录制的多个数据集，但尚无专门针对"跨设备迁移性能"的评估模式。

\textbf{OpenNeuro}\citep{Markiewicz2021OpenNeuro}：专门托管BIDS格式神经影像数据的平台，目前已有60+个EEG数据集，包括LEMON（ds000221）、HBN发布版及数百个任务态数据集，平台统一的机器可读元数据格式极大便利了自动化跨设备融合管线的开发。

\textbf{PhysioNet}：医疗级生理信号数据库\citep{PhysioNetMotor}，涵盖EEGMMIDB、CHB-MIT、Sleep-EDF、SHHS\citep{SHHS2000}等临床和睡眠数据集。

\textbf{BNCI Horizon 2020}（bnci-horizon-2020.eu）：包含20+个欧洲BCI研究数据集，元数据记录完整（含设备型号、参考方式、硬件滤波），是MI、P300和ERP数据集的高质量来源，包括MoreGrasp多电极类型数据集\citep{Schwarz2020MoreGrasp}。

\textbf{中国BCI数据集}：清华大学SSVEP系列（Benchmark SSVEP\citep{WangBenchmarkSSVEP2016}、BETA\citep{BETA2020}，共205名被试，Neuroscan SynAmps2）和上海交大SEED系列\citep{Zheng2015SEED}（ESI NeuroScan）构成了中国特色的跨设备融合研究资源。清华SSVEP语料库是目前最大的单设备SSVEP数据库，可作为跨设备SSVEP融合的标准参考域。

\subsection{数据集元数据质量与BIDS兼容性挑战}

跨设备融合研究的工程可落地性高度依赖数据集元数据的完整性。对17个候选数据集的系统评估（按设备型号、固件/软件版本、硬件滤波、参考方案、电极坐标5个维度评分，满分5分）揭示了严峻现状：

\begin{itemize}
  \item \textbf{仅4/17个数据集达到4分以上}，大多数存在1--3项关键元数据缺失；
  \item \textbf{7个数据集设备型号未知}，直接阻碍了设备效应的可解释性分析；
  \item \textbf{电极坐标是最频繁缺失的字段}——仅少数数据集提供逐被试数字化坐标；
  \item 经典基准（BCI Competition IV\citep{Tangermann2012BCIcomp}、Sleep-EDF）得分最低（0--2分），需研究者在融合前补充"最小元数据"。
\end{itemize}

这一元数据质量危机与第\ref{sec:metadata}章BIDS-EEG规范的核心动机直接对应。BIDS-EEG的强制性字段（\texttt{Manufacturer}、\texttt{ManufacturersModelName}、\texttt{EEGReference}、\texttt{HardwareFilters}）若能在数据发布时被严格填写，将极大降低后续跨设备融合的工程门槛。目前，OpenNeuro平台的60+ EEG数据集已全部通过BIDS验证器，代表了元数据标准化的最高实践水准。

% ============================================================
\section{研究空白与未来方向}
\label{sec:gaps}

基于对现有文献的系统梳理，本章归纳了跨设备EEG融合研究中五个关键的研究空白，并提出相应的技术路线图。

\subsection{五大关键研究空白}

\subsubsection{缺乏"traveling-subject"EEG数据集}

在fMRI领域，让同一批被试在多台扫描仪上采集是量化设备效应和验证协调方法的金标准。而在EEG领域，\textbf{至今没有等价的"旅行被试"数据集}——即同一批被试在5种以上设备上执行相同实验范式的受控多设备采集。D'Amico等（2026）的工作仅覆盖两种设备\citep{DAmico2026}，Lee等（2026）的数据集涵盖多种设备但任务设计较为简单\citep{Lee2026dataset}。创建一个50+被试、在5--8种代表性设备（涵盖科研级、半专业级和消费级）上执行相同范式（静息态、P300、SSVEP、运动想象）的多设备数据集，将是里程碑级别的贡献，且在技术上完全可行。

\subsubsection{EEG协调方法比fMRI滞后约八年}

fMRI领域已建立了成熟的多站点协调方法体系：ComBat（2017--2018）\citep{Fortin2017DTI, Fortin2018cortical}、旅行被试方法、完善的评估框架等。EEG的首个ComBat系统应用出现在2024年\citep{Jaramillo2024ComBat}，且至今缺乏共识性的协调框架。

EEG面临比fMRI更严峻的组合复杂性——数十家设备厂商（vs. 三大MRI供应商），1至256个通道（vs. 标准化体素网格），以及更多的参考方案变种。但神经影像协调方法的成熟体系提供了清晰的蓝图，为EEG领域"弯道超车"提供了理论指导。

\subsubsection{基础模型未显式建模设备效应}

当前的EEG基础模型（LaBraM\citep{Jiang2024LaBraM}、CBraMod\citep{Wang2025CBraMod}、REVE\citep{ElOuahidi2025REVE}）通过灵活的分词机制处理通道配置差异，但将设备身份视为隐式因素。没有模型显式建模或去除设备特异性的噪声特性、放大器传递函数或电极阻抗特征。

\textbf{设备条件化训练}（device-conditioned training）或\textbf{设备对抗训练}（device-adversarial training）是一个尚未系统探索但具有明确理论动机的扩展方向：在预训练阶段加入设备类型嵌入或设备分类对抗分支，迫使模型学习到真正设备不变的神经生理表征。

\subsubsection{缺乏标准化的跨设备评估基准}

现有BCI竞赛（BCI Competition IV\citep{Tangermann2012BCIcomp}、OpenBMI\citep{Lee2019OpenBMI}）和基准平台（MOABB\citep{Chevallier2024MOABB}）在每个数据集中使用单一设备，没有专门的"跨设备"评估模式来隔离设备迁移性能。建立类似GLUE（自然语言处理）的跨设备EEG评估基准，将大幅加速该领域的系统性进展。可行的起点是第\ref{sec:dataset_pairing}节整理的推荐配对方案——MI四厂商组合（EEGMMIDB + Lee2019\_MI + Cho2017 + Dreyer2023，共302名被试）、MAMEM1/MAMEM3 SSVEP配对（研究级vs.消费级）以及情绪识别三层设备梯度（SEED + DEAP + DREAMER）已构成了具备足够多样性的跨设备评估矩阵。

\subsubsection{消费级到研究级的桥梁缺失}

已有超过916项研究使用了消费级EEG设备\citep{Sabio2024scoping}，其中Emotiv占67.7\%，NeuroSky占24.6\%。然而，目前没有经过充分验证的方法能将消费级（1--16通道、干电极、AC耦合）与研究级（64--256通道、湿电极、DC耦合）数据进行有效融合。这一空白阻碍了最具影响力的应用：利用海量消费级数据改进研究和临床模型。

\subsection{可行技术路线图}

基于上述分析，本节提出分阶段推进的可行技术路线图：

\textbf{第一阶段——数据基础设施与基准建设（第1--6个月）}

\begin{itemize}
  \item 采用BIDS-EEG\citep{Pernet2019BIDS}作为统一数据组织格式；
  \item 构建traveling-subject EEG数据集：50+被试在5--8种设备上执行相同范式；
  \item 实施标准化预处理管线：统一重参考（平均参考或REST）、0.5--45\,Hz带通滤波、250\,Hz重采样、球面样条空间插值、基于ICA的伪迹去除；
  \item 仅此数据集本身即可作为里程碑级独立贡献发表。
\end{itemize}

\textbf{第二阶段——协调方法开发（第4--12个月）}

\begin{itemize}
  \item 利用traveling-subject数据集适配并系统验证ComBat变体\citep{Jaramillo2024ComBat}；
  \item 开发设备对抗训练方法：带设备分类对抗分支的神经网络，强制学习设备不变表征；
  \item 系统量化设备效应，发表跨范式设备-被试-会话方差分解结果。
\end{itemize}

\textbf{第三阶段——显式设备建模的基础模型（第6--18个月）}

\begin{itemize}
  \item 扩展LaBraM/REVE架构，加入设备类型嵌入层和设备对抗预训练目标；
  \item 目标：10000+小时多设备预训练数据；
  \item 创建覆盖异常检测、MI分类、情绪识别和睡眠分期的跨设备评估基准。
\end{itemize}

\textbf{第四阶段——联邦跨设备框架（第12--24个月）}

面向EEG设备异构性的联邦学习\citep{Li2020federated}方案：设计设备特异性的局部特征提取器和共享语义表征层，实现跨医院、跨设备的隐私保护EEG模型训练，充分利用电信/云计算基础设施优势。

\subsection{联邦学习与隐私保护}

EEG数据的隐私敏感性（脑电信号可能包含个人健康信息）使得联邦学习成为跨机构协作训练的重要路径\citep{Li2020federated}。在联邦框架下，各医院和研究机构保留本地数据，仅共享模型梯度或参数更新，从而在保护数据隐私的同时实现跨设备模型的联合训练。

\subsection{跨模态对齐的未来展望}

探索更贴合大脑电生理偶极子动力学的自监督对比学习目标，以及更智能的跨模态（将EEG与文本、音频、fMRI结合）对齐机制，是提升下一代模型跨设备、跨被试知识迁移效率的关键命题。EEG与文本的融合已在情感识别领域初见成效\citep{Niso2022HED}，而EEG-fMRI联合预训练则有望为基础模型注入更丰富的空间分辨率信息。

% ============================================================
\section{结论}
\label{sec:conclusion}

本综述系统梳理了跨设备脑电数据融合领域从硬件基础到算法前沿的全景进展。核心结论如下：

\subsection{分层可行性总结}

\begin{enumerate}[leftmargin=2em]
  \item \textbf{原始信号级直接融合：可行性中等偏低}。设备效应约占总信号方差的9\%\citep{Ratti2017systems}，无法忽视。前提是设备间存在较强可比性，必须统一预处理、严格记录参考与硬件滤波，并最好拥有同受试者同步对照数据作为校准基准\citep{DAmico2026}。

  \item \textbf{特征级融合：可行性较高}。对于稳健的频域特征（alpha峰、频带功率）、ERP窗口特征（P300幅值/潜伏期）和协方差矩阵特征，跨设备融合在ComBat、EA、黎曼对齐等方法的支持下已具备较为成熟的方法论基础\citep{Jaramillo2024ComBat, He2020EA}。

  \item \textbf{模型级迁移/域自适应：可行性最高}。AwAR\citep{Wu2017AwAR}、ALPHA\citep{Liu2022ALPHA}、SDDA\citep{Liu2025SDDA}已分别在跨头戴ERP、湿/干电极SSVEP、异构通道布局条件下给出了明确的正向结果。EEG基础模型（REVE\citep{ElOuahidi2025REVE}、LaBraM\citep{Jiang2024LaBraM}等）通过灵活分词和空间位置编码进一步降低了跨设备处理的门槛。
\end{enumerate}

\subsection{当前技术成熟度判断}

计算工具已就绪（Transformer架构、自监督学习、灵活的位置编码），数据基础设施日趋成熟（BIDS-EEG\citep{Pernet2019BIDS}、MOABB\citep{Chevallier2024MOABB}、OpenNeuro\citep{Markiewicz2021OpenNeuro}），科学动机也十分明确。然而，该领域仍面临三个核心未解问题：

\begin{itemize}
  \item 没有traveling-subject基准数据集来系统量化设备效应并验证协调方法；
  \item 基础模型尚未显式建模设备效应；
  \item 缺乏跨设备专用的标准化评估基准。
\end{itemize}

\subsection{战略建议}

综合上述分析，本文提出以下战略建议：

\begin{enumerate}[leftmargin=2em]
  \item \textbf{优先建立traveling-subject数据集}：这是验证所有后续协调方法的基础，也是该领域最迫切的短板。
  \item \textbf{分阶段推进融合策略}：先建立特征级基线，再做域自适应/迁移基线，最后尝试原始信号级统一建模。
  \item \textbf{评估报告双端指标}：同时报告"无自适应下界"和"少量校准/域适应上界"，避免仅报告最优单一结果。
  \item \textbf{拥抱BIDS-EEG标准}：强制记录设备型号、固件版本、采样率、参考方式、硬件滤波、电极类型、通道坐标与同步方式，这是跨设备研究可复现性的必要条件。
\end{enumerate}

跨设备EEG数据融合正从理论上的不可能转变为实践中的现实。随着EEG基础模型的不断成熟和更多traveling-subject数据集的构建，设备无关的通用脑电模型——能够无缝处理从消费级可穿戴设备到高端科研系统的所有EEG数据——正在成为可实现的目标。这不仅将深刻改变认知神经科学和临床神经病学的研究范式，也将为脑机接口技术的大规模工程化应用开辟崭新的可能性。

% ============================================================
\bibliography{references}

\end{document}
